\documentclass[11pt,a4paper]{article}
\usepackage[utf8]{inputenc}
\usepackage[T2A]{fontenc}
\usepackage[russian,english]{babel}
\usepackage{amsmath,amssymb,amsthm}
\usepackage{mathtools}
\usepackage{geometry}
\geometry{margin=2.5cm}
\usepackage{hyperref}
\hypersetup{colorlinks=true,linkcolor=blue!70!black,urlcolor=blue!60!black}
\usepackage{graphicx}
\usepackage{booktabs}
\usepackage{xcolor}
\usepackage{caption}

\title{Чувствительные показатели здоровья почвы: обзор современных подходов}
\author{Обзорная статья}
\date{\today}

\begin{document}

\maketitle

\begin{abstract}
Здоровье почвы --- физическое, химическое и биологическое состояние почвы,
определяющее её способность функционировать как живая система и обеспечивать
экосистемные услуги. Для мониторинга и управления почвами необходимы
индикаторы, которые одновременно чувствительны к агротехническим практикам и
достаточно стабильны, чтобы отражать долгосрочные изменения. В настоящем
обзоре на основе четырёх рецензируемых работ (Feeney et al., 2024; Bai et al.,
2018; Bongiorno et al., 2019; Cusset et al., 2024) рассмотрены три группы
показателей: физико-химические (органический углерод и его нормированные
формы), лабильные фракции углерода и микробиологические индикаторы. Показано,
что наиболее чувствительными и информативными являются лабильные фракции
углерода (в первую очередь POXC) и микробиологические индикаторы (MBC, PLFA,
ферментативная активность), тогда как нормировка содержания углерода по
содержанию глины (SOC/clay) уступает альтернативному подходу O/T SOC.
\end{abstract}

\begin{figure}[h]
  \centering
  \includegraphics[width=0.95\linewidth]{figures/graphical_abstract.pdf}
  \caption{Графический абстракт: связь агротехнических практик, групп
  индикаторов и их чувствительности. Наиболее чувствительные и ранние
  индикаторы --- POXC, MBC, PLFA, ферменты.}
  \label{fig:abstract}
\end{figure}

\section{Введение}
\label{sec:intro}

Деградация земель --- одна из главных экологических проблем человечества.
Европейская обсерватория почв (EUSO) оценивает, что более 60\,\% почв Европы
являются «нездоровыми» (Panagos et al., 2024). В ответ на это Европейский союз
в июле 2023 года предложил Закон о мониторинге почв (Soil Monitoring Law),
направленный на достижение здоровых почв к 2050 году. Улучшение здоровья почв
рассматривается как необходимое условие достижения ряда Целей устойчивого
развития ООН (SDG 2, 6, 13, 15).

Ключевая методологическая проблема --- выбор индикаторов, по которым можно
надёжно оценивать состояние почвы и его изменения. Согласно критерию
Islam \& Weil (2000), лучшие индикаторы --- это свойства, которые не настолько
стабильны, чтобы быть нечувствительными к управлению, но и не настолько
изменчивы, чтобы не давать представления о долгосрочных изменениях. Именно
этот баланс «чувствительности» и «стабильности» лежит в основе выбора
показателей здоровья почвы.

\section{Физико-химические индикаторы: органический углерод и его нормировка}
\label{sec:physchem}

Органический углерод почвы (SOC) --- один из наиболее широко используемых
показателей здоровья почвы наряду с pH и доступными P и K (Bünemann et al.,
2018). SOC влияет на агрегацию, водоудержание, круговорот питательных веществ
и служит средой обитания для биоразнообразия. Однако абсолютное содержание
SOC чувствительно к множеству факторов (климат, тип почвы, землепользование),
что затрудняет определение целевых значений.

\subsection{Отношение SOC/clay}

Поскольку связь с тонкими частицами (глина, ил) --- ключевой механизм
накопления SOC (Hassink, 1997), было предложено нормировать содержание
углерода по содержанию глины. Анализ швейцарских пахотных почв дал три порога
отношения SOC/clay: 1/13, 1/10 и 1/8, разделяющие классы «Деградированная»,
«Умеренная», «Хорошая» и «Очень хорошая» почва (Johannes et al., 2017).
Этот показатель был включён в проект Закона о мониторинге почв ЕС.

\subsection{Индикатор O/T SOC}

Feeney et al. (2024) показали, что нормировка только по глине даёт
вводящие в заблуждение оценки структуры почвы и потенциала накопления
углерода. Авторы развили подход Poeplau \& Don (2023) и предложили индикатор
O/T SOC --- отношение наблюдаемого содержания SOC к «типичному» для данной
педо-климатической зоны (PCZ). Типичное значение --- средняя концентрация SOC
в зоне, учитывающей не только текстуру, но и землепользование и климат.

На основе данных LUCAS (крупнейшая европейская программа мониторинга почв,
2009--2018, $n = 19\,855$ минеральных верхних горизонтов 0--20~см)
регрессионными деревьями выделены 12 педо-климатических зон с типичными
значениями SOC 5,99--39,65~г/кг. Классы индекса установлены по квартилям
распределения O/T SOC в каждой зоне: «Низкий» ($<$ 25-го перцентиля),
«Промежуточный» (25--50), «Высокий» (50--75), «Очень высокий» ($>$ 75).

По сравнению с SOC/clay индикатор O/T SOC:
\begin{itemize}
  \item менее чувствителен к содержанию глины, землепользованию и климату;
  \item менее географически смещён;
  \item лучше отражает различия в пористости и запасах SOC.
\end{itemize}

Эти закономерности подтвердились на двух независимых наборах данных, а классы
O/T SOC оказались достаточно чувствительными, чтобы отражать различия в
управлении на нескольких долгосрочных полевых экспериментах.

\section{Лабильные фракции углерода}
\label{sec:labile}

Изменения, вызванные агротехническими практиками, часто трудно обнаружить по
общему содержанию SOC (Haynes, 2005). Поэтому более информативным считается
измерение быстро меняющихся пулов --- лабильного углерода, который напрямую
доступен для микробной активности и служит основным источником энергии для
микроорганизмов.

Bongiorno et al. (2019) сравнили пять лабильных фракций углерода в 10
европейских долгосрочных полевых экспериментах (LTE):
\begin{itemize}
  \item растворённый органический углерод (DOC);
  \item гидрофильный DOC (Hy-DOC);
  \item углерод, окисляемый перманганатом (POXC, «активный углерод»);
  \item углерод, экстрагируемый горячей водой (HWEC);
  \item углерод органического вещества в виде частиц (POMC).
\end{itemize}

Результаты показали, что сокращённая обработка почвы и высокий внос
органического вещества повышают концентрации лабильных фракций по сравнению
с традиционной обработкой и низким вносом соответственно. Наиболее
чувствительными к обоим факторам (обработка и удобрение) оказались POXC и
POMC. При этом POXC был фракцией, наиболее положительно коррелирующей с
химическими (общий органический углерод, общий азот, ёмкость катионного
обмена), физическими (водоустойчивые агрегаты, влагоудерживающая способность,
плотность) и биологическими (микробная биомасса C и N, дыхание почвы)
показателями.

Авторы заключают, что POXC --- лабильная фракция, чувствительная к управлению
и наиболее информативная о состоянии органического вещества, питательных
веществ, структуры почвы и микробных пулов. Измерение POXC относительно
дёшево, быстро и просто, что делает его рекомендуемым показателем в схемах
оценки качества почвы.

\section{Микробиологические индикаторы}
\label{sec:microbial}

Рутинные тесты почв традиционно опираются на физико-химические компоненты.
Однако микробиологические индикаторы могут быть более чувствительными и
ранними сигналами изменений. Cusset et al. (2024) провели мета-анализ
166 статей, оценив чувствительность активных пулов органического вещества и
микробиологических индикаторов к шести практикам: внос экзогенного
органического вещества (EOM), сохранение растительных остатков (RR), введение
покровных культур (CC), севооборот/диверсификация (CR), сокращение обработки
(TR) и переход к инновационным/органическим системам (IOS). Применялись модели
со случайными эффектами на основе натурального логарифма отношения отклика.

Ключевые результаты:
\begin{itemize}
  \item Климатические факторы значимо не влияли на отклик индикаторов, тогда
  как содержание глины, органического вещества, pH, длительность опыта и
  глубина отбора проб значимо модулировали отклик.
  \item К EOM, RR, CC и TR были чувствительны как физико-химические (SOC,
  общий азот STN), так и микробиологические (MBC, общие PLFA,
  $\beta$-глюкозидаза B-GLU) индикаторы. Например, CC значимо повышала SOC
  (+7\,\%), STN (+8\,\%), MBC (+14\,\%), общие PLFA (+19\,\%) и B-GLU
  (+11\,\%).
  \item CR и IOS не влияли на SOC и STN, но влияли на микробиологические
  индикаторы: MBC вырос на 7--34\,\%, PLFA на 8--24\,\%, а B-GLU --- на
  58\,\% при IOS.
  \item Микробиологические индикаторы (MBC, PLFA) и ферментативная активность
  (FDA, B-GLU, URE) показали большую чувствительность к широкому спектру
  практик независимо от возраста практики, что указывает на их потенциал как
  более отзывчивых, ранних и универсальных индикаторов.
\end{itemize}

Помимо классических микробиологических показателей (биомасса, дыхание,
ферменты), в качестве чувствительных био-индикаторов рассматриваются
молекулярные маркеры. Обзор Bhaduri (2022), обобщивший 250 исследований
по восстановлению экосистем и устойчивости почв, выделяет набор потенциальных
био-индикаторов: микробную биомассу, дыхание, ферментативную активность,
молекулярные генные маркеры, микробные метаболиты и анализ микробных
сообществ, которые отзывчивы к широкому кругу экосистемных функций в
сельскохозяйственных, горнодобывающих, загрязнённых тяжёлыми металлами,
пустынных, радиоактивно и пестицидно загрязнённых и заболоченных почвах
(Bhaduri 2022, стр.~1). Авторы подчёркивают, что определение
эффективных био-индикаторов затруднено из-за их временного и пространственного
разрешения на уровне экосистем (Bhaduri 2022, стр.~1).

Дальнейшим шагом в развитии микробиологических индикаторов стало использование
функциональных генов. Jia (2025) в 10-летнем полевом эксперименте с
удобрениями (пшеница--кукуруза, Северо-Китайская равнина) количественно
определили 17 функциональных генов циклов углерода (cbbL, GH31), азота (nifH,
ureC, chiA, A-amoA, B-amoA, narG, nirK, nirS, norB, nosZ) и фосфора (gltA, bpp,
phoD, phoC, pqqC). Гены показали высокие коэффициенты вариации и сильную
чувствительность к вариантам удобрения при низкой изменчивости внутри одной
обработки; особенно тесно с прокси-показателями циклов C, N и P коррелировали
GH31, cbbL, B-amoA, chiA, phoC и phoD (Jia 2025, стр.~1). Органические
удобрения усиливали численность генов, связанных с углеродом и питательными
веществами, что положительно сказывалось на урожайности кукурузы. Авторы
заключают, что функциональные гены чувствительны к органическим вносам и дают
более детальное и механистическое понимание почвенных процессов, чем обычные
индикаторы (Jia 2025, стр.~1).

Микробные сообщества как целое также могут служить индикаторами. Trivedi
(2016) провели мета-анализ данных из 102 рецензируемых публикаций, оценив, как
свойства, напрямую связанные с питательным статусом почвы (общий C и N,
отношение C:N), первичная продуктивность (NPP) и микробное разнообразие и
состав (относительная численность основных бактериальных фил по данным
секвенирования) реагируют на сельскохозяйственное управление в основных
биомах Земли (аридном, континентальном, умеренном и тропическом). Были
обнаружены устойчивые статистические тренды в относительной численности
нескольких бактериальных фил: Actinobacteria и Chloroflexi --- в агросистемах,
Acidobacteria, Proteobacteria и Cyanobacteria --- в природных системах, причём
эти тренды коррелировали со многими свойствами почвы (Trivedi 2016,
стр.~1). Однако основные эффекты сельского хозяйства на свойства почвы и
продуктивность зависели от биома (Trivedi 2016, стр.~1).

\section{Влияние агротехнических практик на индикаторы}
\label{sec:practices}

Bai et al. (2018) обобщили данные долгосрочных полевых экспериментов Европы и
Китая (30 LTE из 13 площадок проекта iSQAPER, дополненные 42 LTE из Китая и
402 наблюдениями из литературы; рассматривались опыты длительностью не менее
5 лет). Сравнивались четыре пары практик (органическое вещество vs без него,
no-tillage vs традиционная обработка, севооборот vs монокультура, органическое
земледелие vs традиционное) по пяти индикаторам: содержание органического
вещества (SOM), pH, агрегатная стабильность, численность дождевых червей и
урожайность. Относительные эффекты анализировались через отношение отклика
(RR).

Основные выводы:
\begin{itemize}
  \item Внос органического вещества благоприятно влиял на все индикаторы;
  наиболее сильный эффект --- на численность дождевых червей, затем на
  урожайность, SOM и агрегатную стабильность. Влияние на pH зависело от типа
  почвы (положительно для кислых почв).
  \item No-tillage повышал агрегатную стабильность и содержание SOM в верхних
  горизонтах, но урожайность при этом несколько снижалась. Численность червей
  росла, но не там, где применялись гербициды/пестициды.
  \item Севооборот положительно влиял на SOM и урожайность; ротация с травами
  (ley) очень положительно влияла на численность червей.
  \item Органическое земледелие повышало численность червей, агрегатную
  стабильность и содержание SOM, но урожайность снижалась.
\end{itemize}

\section{Синтез: какие индикаторы наиболее чувствительны}
\label{sec:synthesis}

Сопоставление четырёх работ позволяет выделить общую картину.

\textbf{Физико-химические индикаторы} (SOC, STN) надёжно отражают изменения
органического статуса почвы под действием практик, связанных с вносом
углерода (EOM, RR, CC, TR), но менее чувствительны к диверсификации культур и
инновационным системам в краткосрочных исследованиях (Cusset et al., 2024).
Нормировка SOC по глине (SOC/clay) имеет ограничения; более робастным
подходом является O/T SOC, учитывающий педо-климатический контекст
(Feeney et al., 2024).

\textbf{Лабильные фракции углерода}, особенно POXC и POMC, наиболее
чувствительны к обработке почвы и вносу органического вещества, а POXC
дополнительно коррелирует с широким набором физико-химических и биологических
показателей (Bongiorno et al., 2019).

\textbf{Микробиологические индикаторы} (MBC, PLFA, ферментативная активность)
проявляют наибольшую чувствительность к широкому спектру практик, включая
те, которые не меняют SOC/STN в краткосрочной перспективе (CR, IOS), и потому
могут служить ранними индикаторами перехода к устойчивому земледелию
(Cusset et al., 2024).

Сопоставление с более широкой литературой подтверждает эту картину. Cardoso
et al. (2013) подчёркивают, что биологические индикаторы реагируют быстрее
физико-химических, и рекомендуют системный подход, сочетающий физические,
химические и биологические свойства, как более надёжный для оценки здоровья
почвы (Cardoso 2013, стр.~1). Trivedi (2016) показывают, что
реакция микробных сообществ на тип растительности предсказуема, что может
быть использовано для разработки нового набора индикаторов первичной
продуктивности и здоровья почвы (Trivedi 2016, стр.~1). Наконец,
функциональные гены микробного цикла C, N и P, по данным Jia (2025),
чувствительны к органическим вносам и дают более механистическое понимание
почвенных процессов, чем обычные индикаторы, что делает их перспективными
для продвинутых схем оценки здоровья почвы (Jia 2025, стр.~1).

Таким образом, для комплексной оценки здоровья почвы целесообразно сочетать
все три группы: физико-химические показатели (с нормировкой по
педо-климатическому контексту), лабильные фракции углерода (POXC) и
микробиологические индикаторы (MBC, PLFA, ферменты). Именно микробиологические
и лабильные показатели оказываются наиболее чувствительными и ранними, что
делает их приоритетными для включения в рутинные лабораторные тесты.

\section{Сводная таблица индикаторов}
\label{sec:table}

В таблице~\ref{tab:indicators} сведены рассмотренные индикаторы, их
чувствительность к агротехническим практикам, информативность и практичность
(стоимость/скорость измерения). Оценки качественные, основаны на выводах
четырёх анализируемых работ.

\begin{table}[h]
  \centering
  \caption{Сводка индикаторов здоровья почвы.}
  \label{tab:indicators}
  \small
  \begin{tabular}{lcccc}
    \toprule
    Индикатор & Группа & Чувств. & Информ. & Практич. \\
    \midrule
    SOC / STN & физико-хим. & средняя & высокая & высокая \\
    SOC/clay & физико-хим. & низкая & средняя & высокая \\
    O/T SOC & физико-хим. & средняя--высокая & высокая & средняя \\
    POXC & лабильный C & высокая & высокая & высокая \\
    POMC & лабильный C & высокая & средняя & средняя \\
    DOC / HWEC & лабильный C & средняя & средняя & средняя \\
    MBC & микробиол. & высокая & высокая & средняя \\
    PLFA & микробиол. & высокая & высокая & низкая \\
    Ферменты (B-GLU, FDA, URE) & микробиол. & высокая & высокая & средняя \\
    Дождевые черви & биологич. & высокая & средняя & низкая \\
    Агрегатная стабильность & физич. & средняя & высокая & средняя \\
    \bottomrule
  \end{tabular}
\end{table}

\textbf{Вывод по таблице.} Наиболее чувствительные и одновременно практичные
индикаторы --- POXC (лабильный углерод) и MBC (микробиологическая биомасса).
PLFA и ферментативная активность максимально чувствительны, но требуют более
сложных лабораторных протоколов. Физико-химические показатели (SOC, STN)
остаются базовыми, а их нормировка по педо-климатическому контексту (O/T SOC)
повышает надёжность оценки.

\section{Ограничения существующих индикаторов}
\label{sec:limitations}

Несмотря на высокую чувствительность рассмотренных показателей, каждый из
них имеет существенные ограничения, которые необходимо учитывать при
интерпретации результатов мониторинга.

\textbf{Нормировка SOC по педо-климатическому контексту применима не ко всем
почвам.} Индикатор O/T SOC разработан только для минеральных верхних
горизонтов: для органических почв ($\geq$ 120~г/кг SOC) изменчивость
концентраций значительно выше, и неясно, способны ли классы O/T SOC разделять
различающиеся градации структуры органических почв; кроме того, для
подпочвенных горизонтов, вероятно, потребуется иной набор «типичных»
концентраций SOC (Feeney et al., 2024, стр.~10).

\textbf{Пороги и классы O/T SOC зависят от данных и цифровых карт почв.}
Распределение педо-климатических зон (и вместе с ними порогов O/T SOC)
меняется со временем по мере изменения земного покрова, а предсказанные классы
чувствительны к точности цифровых карт почв, которые часто плохо предсказывают
экстремальные значения; в результате пользователь может получить вводящие в
заблуждение оценки, отражающие слабости самих карт, а не реальное состояние
почвы (Feeney et al., 2024, стр.~9).

\textbf{Агрегирование данных исключает анализ динамики во времени.} Для
построения O/T SOC данные всех трёх доступных съёмок LUCAS были объединены,
что не позволяет проводить анализ изменений во времени на основе самого LUCAS;
кроме того, требуются дополнительные данные, чтобы подтвердить, что рост
O/T SOC действительно ведёт к увеличению пористости и запасов SOC
(Feeney et al., 2024, стр.~10).

\textbf{Лабильные фракции углерода чувствительны к условиям отбора проб.}
Некоторые исследования не обнаружили эффекта обработки и удобрения на
лабильные фракции, что может объясняться свойствами почвы, неоднородным
распределением растительных и микробных остатков, типом и количеством
вносимого органического вещества, условиями окружающей среды и временем
отбора проб; кроме того, о химическом составе и сезонной динамике POXC
известно мало, а его чувствительность к другим практикам (например, покровным
культурам) требует дальнейшей валидации (Bongiorno et al., 2019, стр.~10).

\textbf{Микробиологические индикаторы дороже и сложнее в рутинном анализе.}
По сравнению с физико-химическими показателями анализ активных пулов
органического вещества и микробиологических индикаторов требует больше времени
на отбор проб и отправку в лабораторию, более высокой технической
квалификации, большей длительности и стоимости; для включения в рутинные тесты
индикаторы должны анализироваться по стандартизированным протоколам и не
требовать высокой технической экспертизы (Cusset et al., 2024, стр.~2).

\textbf{Методологические смещения в мета-анализах.} В мета-анализе Cusset et
al. (2024) использовалась стратегия отбора проб на фиксированную глубину,
несмотря на возможные смещения, связанные с различиями в плотности почвы между
обработками, --- выбор, сделанный с учётом потенциального использования этих
индикаторов в рутинных тестах (Cusset et al., 2024, стр.~2).

\textbf{Неравномерная обеспеченность данными в обзорных исследованиях.} В
обзоре Bai et al. (2018) данных по урожайности, содержанию органического
вещества и pH было значительно больше, чем по численности дождевых червей,
что представляет собой известное ограничение такого рода дескриптивных
исследований; для ограничения влияния возможных выбросов использовались
медианы вместо средних (Bai et al., 2018, стр.~3).

\textbf{Эффекты микробных индикаторов зависят от биома.} Мета-анализ Trivedi
(2016) показал, что основные эффекты сельского хозяйства на свойства почвы и
первичную продуктивность были зависимы от биома: аридные, континентальные,
умеренные и тропические экосистемы могут реагировать на сельскохозяйственную
эксплуатацию по-разному, что ограничивает универсальность микробных
индикаторов при переносе между регионами (Trivedi 2016, стр.~1).

\textbf{Функциональные гены требуют сложных протоколов.} Хотя функциональные
гены чувствительны к органическим вносам, их количественное определение
(17 генов циклов C, N и P) требует молекулярно-биологических методов и
специализированной лабораторной инфраструктуры, что ограничивает их применение
в рутинных тестах по сравнению с более простыми показателями (Jia 2025,
стр.~1).

\textbf{Определение эффективных био-индикаторов затруднено их разрешением.}
Как отмечают Bhaduri (2022), выбор эффективных био-индикаторов здоровья
почвы остаётся сложной задачей из-за их временного и пространственного
разрешения на уровне экосистем: один и тот же показатель может быть
чувствительным в одних условиях и малоинформативным в других (Bhaduri 2022,
стр.~1).

\section{Заключение}
\label{sec:conclusion}

Обзор показывает, что выбор чувствительных индикаторов здоровья почвы ---
ключевая задача для мониторинга и управления. Наиболее чувствительными и
информативными являются лабильные фракции углерода (POXC) и микробиологические
индикаторы (MBC, PLFA, ферментативная активность), тогда как традиционные
физико-химические показатели дополняют их, особенно при нормировке по
педо-климатическому контексту (O/T SOC). Дальнейшие исследования должны быть
направлены на стандартизацию протоколов измерения и валидацию индикаторов в
различных педо-климатических условиях.

\begin{thebibliography}{9}
\bibitem{feeney2024}
Feeney, C.J., Bentley, L., De Rosa, D., Panagos, P., Emmett, B.A., Thomas, A.,
Robinson, D.A. (2024). Benchmarking soil organic carbon (SOC) concentration
provides more robust soil health assessment than the SOC/clay ratio at European
scale. \emph{Science of the Total Environment}, 951, 175642.

\bibitem{bai2018}
Bai, Z., Caspari, T., Ruiperez Gonzalez, M., Batjes, N.H., Mäder, P.,
Bünemann, E.K., et al. (2018). Effects of agricultural management practices on
soil quality: A review of long-term experiments for Europe and China.
\emph{Agriculture, Ecosystems and Environment}, 265, 1--7.

\bibitem{bongiorno2019}
Bongiorno, G., Bünemann, E.K., Oguejiofor, C.U., Meier, J., Gort, G., Comans, R.,
Mäder, P., Brussaard, L., de Goede, R. (2019). Sensitivity of labile carbon
fractions to tillage and organic matter management and their potential as
comprehensive soil quality indicators across pedoclimatic conditions in Europe.
\emph{Ecological Indicators}, 99, 38--50.

\bibitem{cusset2024}
Cusset, E., Bennegadi-Laurent, N., Recous, S., Bernard, P.-Y., Perrin, A.-S.,
Tscheiller, R., Trinsoutrot-Gattin, I., Riah-Anglet, W. (2024). Which soil
microbial indicators should be included in routine laboratory tests to support
the transition to sustainable management of arable farming systems?
A meta-analysis. \emph{Ecological Indicators}, 167, 112706.

\bibitem{bhaduri2022}
Bhaduri, D., Sihi, D., Bhowmik, A., Verma, B.C., Munda, S., Dari, B. (2022).
A review on effective soil health bio-indicators for ecosystem restoration and
sustainability. \emph{Frontiers in Microbiology}, 13, 938481.

\bibitem{cardoso2013}
Cardoso, E.J.B.N., Vasconcellos, R.L.F., Bini, D., Miyauchi, M.Y.H.,
Santos, C.A., Alves, P.R.L., de Paula, A.M., Nakatani, A.S., Pereira, J.M.,
Nogueira, M.A. (2013). Soil health: looking for suitable indicators. What
should be considered to assess the effects of use and management on soil
health? \emph{Scientia Agricola}, 70(4), 274--289.

\bibitem{jia2025}
Jia, J., de Goede, R., Li, Y., Zhang, J., Wang, G., Zhang, J., Creamer, R.
(2025). Unlocking soil health: Are microbial functional genes effective
indicators? \emph{Soil Biology and Biochemistry}, 204, 109768.

\bibitem{trivedi2016}
Trivedi, P., Delgado-Baquerizo, M., Anderson, I.C., Singh, B.K. (2016).
Response of soil properties and microbial communities to agriculture:
Implications for primary productivity and soil health indicators.
\emph{Frontiers in Plant Science}, 7, 990.

\bibitem{panagos2024}
Panagos, P., Broothaerts, N., Ballabio, C., Orgiazzi, A., De Rosa, D.,
Borrelli, P., et al. (2024). How the EU Soil Observatory is providing solid
science for healthy soils. \emph{European Journal of Soil Science}, 75, e13507.

\bibitem{bunemann2018}
Bünemann, E.K., Bongiorno, G., Bai, Z., Creamer, R., De Deyn, G., de Goede, R.,
et al. (2018). Soil quality --- a critical review. \emph{Soil Biology and
Biochemistry}, 120, 105--125.
\end{thebibliography}

\end{document}